\begin{enumerate}
	\item On obtient successivement :
	\begin{itemize}
		\item $1- 2$ donc $-1$

		\item $(-1)\times 5$ donc $-5$

		\item $-5+3\times 1$ donc  $-5 + 3$ soit $-2$
	\end{itemize}
	Le résultat est $- 2$.

	\item On note $x$ le nombre choisi au départ.

	On obtient successivement :
	\begin{itemize}
		\item $x - 2$

		\item $5(x - 2)$ donc $5x - 10$

		\item $5x-10+3x$ donc $8x - 10$
	\end{itemize}

	Le résultat final est $8x-10$.

	\item
$\begin{array}[t]{@{} r r !{=} l}
\text{On résout l'équation: } & 8x-10 & 0\\
 \text{équivaut à: } & 8x & 10\\
 \text{équivaut à: } & x & \dfrac{10}{8}\\
 \text{équivaut à: } & x & 1,25
\end{array}$


	Pour que le résultat soit égal à 0, le nombre de départ doit être 1,25.

	\begin{enumerate}
		\item En partant de 1, on obtient successivement :
		\begin{itemize}
\item $4\times 1$ donc $4$
\item $4 - 5$ donc $-1$
\item $2\times (- 1)$ donc $-2$
		\end{itemize}

	Le résultat est $-2$.
		\item On applique le programme à un nombre $x$ quelconque.

	On obtient successivement :
\begin{itemize}
\item $4\times x$ donc $4x$
\item $4x-5$
\item $(4x-5)\times 2$ donc $8x - 10$
		\end{itemize}
On obtient donc le même résultat qu'avec le programme A.
	\end{enumerate}
\end{enumerate}

\needspace{10\baselineskip}
